\SourceLeaf{348}%
negative for concave-convex surfaces. If the patch under consideration consists of parts of different kinds in this respect, further specifications are needed and must be passed over here.

Comparing the areas of corresponding patches of the curved surface and the auxiliary sphere leads to a new concept, much as comparing volume and mass leads to the concept of density. The author calls the value of a fraction the \textit{measure of curvature} at a point on the curved surface. Its denominator is the area of an infinitely small patch of the curved surface at that point. Its numerator is the area of the corresponding patch of the auxiliary sphere, or the total curvature of that element. One can see that, in the author's sense, total curvature and measure of curvature on curved surfaces are analogous, respectively, to amplitude and curvature itself for curved lines. He had reservations about transferring the latter terms to curved surfaces, since they owe their acceptance more to habit than to fitness. The names themselves are less important than having substantive theorems to justify introducing them.

The problem of finding the measure of curvature at each point of a curved surface takes different forms according to how the surface is specified. The simplest is to distinguish points in space by three rectangular coordinates \(x,y.\ z\) and express one coordinate as a function of the other two. This gives the simplest expression for the measure of curvature. It also reveals a remarkable connection between this measure and the curvatures of the curves formed by intersecting the surface with planes normal to it. As is well known, \textsc{Euler} first showed that there are two such cutting planes, themselves at right angles to one another, with the greatest radius of curvature in one and the smallest in the other, or more precisely, with the two extreme curvatures occurring in them. The expression for the measure of curvature just mentioned shows that it equals a fraction with numerator one and denominator the product of the two extreme radii of curvature. The expression for the measure of curvature is less simple when
